\chapter{Estado del arte}
\label{cap:estado-del-arte}

Este cap\'itulo revisa los fundamentos que permiten estudiar la influencia del
codificador visual en una pol\'itica de difusi\'on. La exposici\'on parte del
aprendizaje por imitaci\'on y contin\'ua con los modelos generativos de difusi\'on.
Despu\'es, se describe \textit{Diffusion Policy} y se analizan sus principales
extensiones.

A continuaci\'on, la revisi\'on se centra en las representaciones visuales empleadas
en este trabajo. Se consideran las redes residuales, los transformadores visuales y
el preentrenamiento supervisado, autosupervisado y multimodal. El cap\'itulo concluye
delimitando el hueco experimental que aborda la metodolog\'ia.

\section{Aprendizaje por imitaci\'on en manipulaci\'on rob\'otica}
\label{sec:imitacion}

El aprendizaje por imitaci\'on obtiene una pol\'itica a partir de demostraciones de un
experto. Su formulaci\'on m\'as directa, denominada clonaci\'on del comportamiento,
trata cada par observaci\'on-acci\'on como un ejemplo supervisado. Este planteamiento
evita dise\~nar una funci\'on de recompensa, pero depende de la cobertura y calidad de
las demostraciones \cite{ross2011dagger}, \cite{mandlekar2021robomimic}.

Sin embargo, el entrenamiento supervisado y la ejecuci\'on de la pol\'itica siguen
distribuciones diferentes. Un error modifica el estado visitado y puede llevar al
robot fuera de los datos de entrenamiento. Los errores posteriores se acumulan
porque la pol\'itica carece de ejemplos en esas nuevas regiones. DAgger reduce este
desplazamiento mediante la agregaci\'on iterativa de correcciones del experto
\cite{ross2011dagger}.

Adem\'as, las demostraciones de manipulaci\'on suelen contener varias acciones
v\'alidas para una misma observaci\'on. Una p\'erdida cuadr\'atica aplicada a una
pol\'itica determinista aproxima la media condicional de esas alternativas. Dicha
media puede situarse entre modos v\'alidos y generar una acci\'on que no aparece en
ninguna demostraci\'on \cite{florence2021ibc}.

Para representar esa multimodalidad, la clonaci\'on de comportamiento impl\'icita
define una funci\'on de energ\'ia sobre pares de observaci\'on y acci\'on. La acci\'on se
obtiene minimizando esa energ\'ia durante la inferencia. Este enfoque mejora la
expresividad, aunque requiere muestras negativas y una optimizaci\'on adicional
\cite{florence2021ibc}. Estas limitaciones motivan el uso de modelos generativos
directamente sobre las secuencias de control.

\section{Fundamentos de los modelos de difusi\'on}
\label{sec:fundamentos-difusion}

Los modelos de difusi\'on aprenden a invertir una degradaci\'on gradual de los datos.
Los modelos probabil\'isticos de difusi\'on por eliminaci\'on de ruido (DDPM) definen
un proceso directo que a\~nade ruido gaussiano durante $K$ pasos. Para una muestra
limpia $\mathbf{x}^{0}$, un estado perturbado puede expresarse como
\cite{ho2020ddpm}

\begin{equation}
  q(\mathbf{x}^{k}\mid\mathbf{x}^{0}) =
  \mathcal{N}\!\left(\sqrt{\bar{\alpha}_{k}}\,\mathbf{x}^{0},
  (1-\bar{\alpha}_{k})\mathbf{I}\right).
  \label{eq:proceso-directo}
\end{equation}

En cambio, el proceso inverso parte de ruido y estima sucesivamente una muestra de
la distribuci\'on original. Ho \textit{et al.} \cite{ho2020ddpm} parametrizan una red
que predice el ruido incorporado en cada paso. Su objetivo simplificado es un error
cuadr\'atico entre el ruido real y el estimado, lo que permite un entrenamiento
estable.

Asimismo, los modelos basados en puntuaci\'on aprenden la direcci\'on de mayor
crecimiento de la densidad logar\'itmica. Este campo permite generar muestras
mediante din\'amica de Langevin
\cite{song2019ncsn}. La introducci\'on de varios niveles de ruido facilita la
estimaci\'on en regiones con pocos datos.

Por otra parte, las ecuaciones diferenciales estoc\'asticas (SDE) ofrecen una
formulaci\'on continua com\'un para estos procesos. El proceso inverso depende de la
funci\'on de puntuaci\'on de las distribuciones perturbadas. Esta perspectiva conecta
los DDPM con otros modelos basados en puntuaci\'on y ampl\'ia las alternativas de
muestreo \cite{song2021sde}.

No obstante, el muestreo iterativo aumenta la latencia. Los modelos impl\'icitos de
difusi\'on por eliminaci\'on de ruido (DDIM) conservan el objetivo de entrenamiento de
los DDPM, pero emplean un proceso inverso no markoviano. Esta modificaci\'on reduce el
n\'umero de evaluaciones necesarias durante la generaci\'on \cite{song2021ddim}.

En conjunto, estos trabajos proporcionan tres elementos para el control rob\'otico:
una distribuci\'on generativa expresiva, un objetivo de predicci\'on de ruido y un
procedimiento iterativo de muestreo. \textit{Diffusion Policy} traslada esos
elementos desde la generaci\'on de datos al espacio de acciones.

\section{Diffusion Policy para control visuomotor}
\label{sec:diffusion-policy}

\textit{Diffusion Policy} modela la distribuci\'on condicional
$p(\mathbf{A}_{t}\mid\mathbf{O}_{t})$, donde $\mathbf{O}_{t}$ contiene las
observaciones recientes y $\mathbf{A}_{t}$ representa una secuencia de acciones
futuras \cite{chi2025diffusion}. Durante el entrenamiento, se perturba una secuencia
demostrada y se minimiza la p\'erdida

\begin{equation}
  \mathcal{L}_{\mathrm{DP}} =
  \mathbb{E}_{k,\boldsymbol{\epsilon}}
  \left[\left\lVert
  \boldsymbol{\epsilon} -
  \boldsymbol{\epsilon}_{\theta}
  (\mathbf{O}_{t},\mathbf{A}_{t}^{k},k)
  \right\rVert_{2}^{2}\right].
  \label{eq:perdida-diffusion-policy}
\end{equation}

De este modo, la red $\boldsymbol{\epsilon}_{\theta}$ aprende a eliminar el ruido de
una trayectoria condicionada por la observaci\'on. La formulaci\'on admite varios
modos de acci\'on sin promediarlos en una \'unica salida. Adem\'as, evita estimar la
constante de normalizaci\'on que dificulta el entrenamiento de los modelos de energ\'ia
\cite{chi2025diffusion}.

En cuanto a la arquitectura, el trabajo original propone dos redes para predecir el
ruido. La primera es una red convolucional temporal con estructura U-Net y
modulaci\'on lineal por caracter\'isticas (FiLM). La segunda emplea un transformador
temporal con atenci\'on cruzada. La variante convolucional ofrece mayor estabilidad,
mientras que el transformador responde mejor ante cambios de acci\'on r\'apidos
\cite{chi2025diffusion}.

Adem\'as, la pol\'itica aplica control de horizonte recesivo. A partir de una ventana
de observaciones, predice una secuencia completa y ejecuta solo sus primeras
acciones. Despu\'es, incorpora una nueva observaci\'on y repite la predicci\'on. Esta
estrategia combina coherencia temporal con capacidad de reacci\'on ante desviaciones
\cite{chi2025diffusion}.

En la rama visual, cada imagen se transforma en un vector antes del proceso de
eliminaci\'on de ruido. La implementaci\'on de referencia utiliza una ResNet-18 sin
preentrenamiento. El promedio espacial global se sustituye por un \textit{spatial
softmax} para conservar informaci\'on posicional. Asimismo, la normalizaci\'on por
lotes se reemplaza por normalizaci\'on por grupos para estabilizar el entrenamiento
\cite{chi2025diffusion}.

La influencia de esta rama ya se examina parcialmente en el art\'iculo original. La
\cref{tab:ablation-visual-dp} resume la comparaci\'on realizada sobre la tarea
\textit{Square} de RoboMimic. Se mantienen fijos el conjunto de demostraciones y la
pol\'itica convolucional, pero cambian la arquitectura y la estrategia de adaptaci\'on
del codificador.

\begin{table}[htbp]
  \centering
  \caption{Resultado de la ablaci\'on de codificadores visuales de \textit{Diffusion
  Policy} en la tarea \textit{Square}.}
  \label{tab:ablation-visual-dp}
  \begin{tabular}{lccc}
    \toprule
    Codificador y preentrenamiento & Desde cero & Congelado & Ajuste fino \\
    \midrule
    ResNet-18, ImageNet-21k & \num{0.94} & \num{0.58} & \num{0.92} \\
    ResNet-34, ImageNet-21k & \num{0.92} & \num{0.40} & \num{0.94} \\
    ViT-B/16, CLIP           & \num{0.22} & \num{0.70} & \num{0.98} \\
    \bottomrule
  \end{tabular}
  \fuente{adaptado de \cite[Tabla~5]{chi2025diffusion}}
\end{table}

Como muestra la tabla, ninguna estrategia domina para todas las arquitecturas. El
ajuste fino obtiene el mejor resultado con ViT-B/16, mientras que ResNet-18 desde
cero supera ligeramente a su versi\'on ajustada. En cambio, la congelaci\'on reduce el
rendimiento de ambas redes residuales \cite{chi2025diffusion}.

Por \'ultimo, Push-T permite estudiar estas decisiones en una tarea de contacto
planar. Un efector circular debe empujar una pieza con forma de T hasta una regi\'on
objetivo. La puntuaci\'on depende del solapamiento entre la pieza y dicha regi\'on, por
lo que exige controlar posici\'on y orientaci\'on \cite{florence2021ibc},
\cite{chi2025diffusion}. En la versi\'on f\'isica, la pol\'itica entrenada de extremo a
extremo alcanz\'o un \num{0.95} de \'exito, frente a \num{0.80} con R3M congelado y
\num{0.15} con ImageNet congelado \cite{chi2025diffusion}. Estos resultados confirman
que la representaci\'on visual puede alterar el control aun cuando la red generativa
permanece fija.

\section{Evoluci\'on de las pol\'iticas de difusi\'on}
\label{sec:evolucion-diffusion-policy}

Tras la publicaci\'on de \textit{Diffusion Policy}, la investigaci\'on se ha dividido
en varias l\'ineas. Unos trabajos modifican la representaci\'on sensorial, mientras
que otros act\'uan sobre la jerarqu\'ia, la geometr\'ia o la velocidad de inferencia.
La \cref{tab:evolucion-dp} sintetiza las extensiones m\'as relacionadas con este
estudio.

\begin{table}[htbp]
  \centering
  \small
  \caption{Principales l\'ineas de extensi\'on de \textit{Diffusion Policy}.}
  \label{tab:evolucion-dp}
  \begin{tabular}{>{\raggedright\arraybackslash}p{2.6cm}
                  >{\raggedright\arraybackslash}p{3.2cm}
                  >{\raggedright\arraybackslash}p{3.4cm}
                  >{\raggedright\arraybackslash}p{4.5cm}}
    \toprule
    Trabajo & Componente modificado & Representaci\'on & Contribuci\'on principal \\
    \midrule
    DP3 \cite{ze2024dp3} & Percepci\'on & Nube de puntos & Sustituye la imagen 2D por
    una representaci\'on geom\'etrica 3D compacta. \\
    HDP \cite{ma2024hdp} & Jerarqu\'ia y cinem\'atica & RGB-D y nube de puntos &
    Separa la planificaci\'on de alto nivel del control difusivo local. \\
    Consistency Policy \cite{prasad2024consistency} & Inferencia & Imagen 2D &
    Destila la trayectoria de eliminaci\'on de ruido para reducir la latencia. \\
    Percepci\'on activa \cite{sun2024active} & Espacio de estado y acci\'on & C\'amara
    m\'ovil RGB & Coordina el punto de vista mediante una restricci\'on cinem\'atica. \\
    ET-SEED \cite{tie2025etseed} & Equivarianza geom\'etrica & Nube de puntos &
    Introduce simetr\'ia SE(3) en la generaci\'on de trayectorias. \\
    DINOv3-DP \cite{egbe2025dinov3dp} & Codificador visual & Imagen 2D & Compara
    ResNet-18 y DINOv3 bajo tres estrategias de entrenamiento. \\
    \bottomrule
  \end{tabular}
  \fuente{elaboraci\'on propia}
\end{table}

En primer lugar, DP3, HDP y ET-SEED incorporan estructura geom\'etrica mediante
nubes de puntos o restricciones cinem\'aticas \cite{ze2024dp3}, \cite{ma2024hdp},
\cite{tie2025etseed}. Estas propuestas mejoran la generalizaci\'on espacial, pero
cambian simult\'aneamente la modalidad y la arquitectura perceptiva. Por ello, sus
resultados no permiten aislar el efecto de un codificador de imagen 2D.

En segundo lugar, \textit{Consistency Policy} reduce el coste del muestreo mediante
destilaci\'on. El trabajo mantiene constante la codificaci\'on de imagen y concentra
la comparaci\'on en la etapa generativa \cite{prasad2024consistency}. De forma
an\'aloga, la percepci\'on activa modifica el espacio de control para coordinar la
c\'amara y el manipulador \cite{sun2024active}.

Finalmente, el estudio DINOv3-DP es el antecedente m\'as pr\'oximo a esta
investigaci\'on. Compara ResNet-18 y un ViT-S/16 autosupervisado en Push-T, Lift, Can
y Square. Adem\'as, eval\'ua entrenamiento desde cero, congelaci\'on y ajuste fino
durante 100 \'epocas \cite{egbe2025dinov3dp}. Sus resultados muestran que la ventaja
del preentrenamiento depende de la tarea y de la estrategia de adaptaci\'on.

\section{Codificadores visuales y preentrenamiento}
\label{sec:codificadores-visuales}

La revisi\'on anterior muestra que la imagen condiciona toda la secuencia de acciones.
Por tanto, el codificador debe conservar la geometr\'ia relevante y descartar
variaciones que no afecten a la tarea. Las familias seleccionadas difieren tanto en
su arquitectura como en la se\~nal utilizada durante el preentrenamiento.

\subsection{Redes residuales y supervisi\'on en ImageNet}
\label{subsec:resnet-imagenet}

Las redes residuales aprenden funciones de correcci\'on respecto a la entrada de cada
bloque. Las conexiones de identidad facilitan la optimizaci\'on de redes profundas y
permiten reutilizar caracter\'isticas en distintas escalas \cite{he2016resnet}. En
particular, ResNet-18 ofrece un compromiso entre capacidad y coste que favorece su
uso en pol\'iticas visuomotoras.

Asimismo, el preentrenamiento supervisado en ImageNet emplea etiquetas de objeto
sobre un conjunto visual amplio \cite{deng2009imagenet}. Estos pesos proporcionan
detectores generales de textura, forma y categor\'ia. Sin embargo, la clasificaci\'on
no exige conservar toda la posici\'on espacial que requiere una tarea de contacto.

Por esta raz\'on, una representaci\'on adecuada para reconocimiento puede no serlo
para control.

\subsection{Transformadores visuales y aprendizaje autosupervisado}
\label{subsec:vit-dinov2}

En cambio, el transformador visual (ViT) divide la imagen en parches y los procesa
como una secuencia de \textit{tokens}. La autoatenci\'on relaciona regiones alejadas
sin depender exclusivamente de la localidad convolucional. Su rendimiento mejora
cuando existe un preentrenamiento a gran escala \cite{dosovitskiy2021vit}.

Sobre esta arquitectura, DINOv2 aprende mediante autodestilaci\'on sin etiquetas. El
m\'etodo combina datos visuales diversos y curados con una estrategia de
entrenamiento autosupervisado. Sus caracter\'isticas transfieren a tareas de imagen y
de nivel de p\'ixel sin ajuste espec\'ifico \cite{oquab2024dinov2}.

En consecuencia, DINOv2 resulta pertinente para manipulaci\'on con pocas
demostraciones. Sus representaciones densas pueden aportar correspondencias
espaciales aprendidas fuera del dominio rob\'otico. Sin embargo, la congelaci\'on
impide adaptar esas caracter\'isticas a la din\'amica y a los contactos de Push-T. La
comparaci\'on con DINOv3 confirma que esta decisi\'on debe evaluarse para cada tarea
\cite{egbe2025dinov3dp}.

\subsection{Preentrenamiento multimodal con CLIP}
\label{subsec:clip}

Por su parte, el preentrenamiento contrastivo de imagen y lenguaje (CLIP) alinea una
imagen con su descripci\'on textual. El modelo original utiliza 400 millones de pares
obtenidos de Internet y transfiere sus representaciones a distintas tareas visuales
\cite{radford2021clip}. La supervisi\'on ling\"u\'istica favorece caracter\'isticas con
contenido sem\'antico general.

Sin embargo, Push-T no requiere interpretar instrucciones ni reconocer categor\'ias
abiertas. La pol\'itica necesita localizar una pieza, estimar su orientaci\'on y
relacionarla con una meta. Por tanto, la sem\'antica de CLIP solo resultar\'a \'util si
su representaci\'on conserva tambi\'en la informaci\'on geom\'etrica necesaria.

Adem\'as, la ablaci\'on original muestra una diferencia marcada entre congelaci\'on y
ajuste fino. El ViT-B/16 preentrenado con CLIP obtiene \num{0.70} cuando permanece
fijo y \num{0.98} cuando se adapta \cite{chi2025diffusion}. Este resultado impide
suponer que un corpus de preentrenamiento mayor garantiza por s\'i solo una pol\'itica
mejor.

\subsection{Representaciones visuales preentrenadas para control}
\label{subsec:representaciones-control}

Frente al preentrenamiento visual gen\'erico, otra l\'inea utiliza datos y objetivos
relacionados con la interacci\'on. R3M aprende de v\'ideo egoc\'entrico mediante
objetivos temporales y alineaci\'on con lenguaje. Como m\'odulo congelado, mejora
varias tareas de manipulaci\'on con pocos datos \cite{nair2022r3m}. No obstante, en
Push-T queda por debajo del codificador entrenado de extremo a extremo
\cite{chi2025diffusion}.

Asimismo, MVP preentrena un ViT mediante modelado enmascarado sobre im\'agenes de
interacci\'on humana. Despu\'es, mantiene fijo el codificador y aprende controladores
mediante aprendizaje por refuerzo. En las ocho tareas de PixMC, esta estrategia
supera al preentrenamiento supervisado en ImageNet en siete casos
\cite{xiao2022mvp}. Sin embargo, el estudio no utiliza demostraciones ni una
pol\'itica de difusi\'on.

Por su parte, Voltron combina reconstrucci\'on visual enmascarada con supervisi\'on
ling\"u\'istica. Su evaluaci\'on sobre cinco problemas rob\'oticos muestra que las
representaciones reconstructivas favorecen detalles espaciales, mientras que los
objetivos contrastivos capturan mejor la sem\'antica. El equilibrio entre ambos
niveles depende de la tarea, y una mayor carga sem\'antica puede perjudicar algunos
problemas de control de bajo nivel \cite{karamcheti2023voltron}.

Finalmente, CortexBench compara CLIP, R3M, MVP y VIP sobre 17 tareas de siete bancos
de pruebas. Ninguna representaci\'on obtiene el mejor resultado en todos los dominios.
VC-1 alcanza el mejor promedio, pero tampoco domina cada tarea; adem\'as, su adaptaci\'on
espec\'ifica produce mejoras sustanciales \cite{majumdar2023vc1}. Por tanto, la escala
del preentrenamiento no determina por s\'i sola la utilidad de un codificador para
control.

\subsection{Estrategias de adaptaci\'on}
\label{subsec:estrategias-adaptacion}

En t\'erminos generales, el entrenamiento desde cero adapta toda la representaci\'on a
la tarea, pero depende exclusivamente de las demostraciones disponibles. Esta opci\'on
reduce el sesgo entre preentrenamiento y control a costa de una mayor necesidad de
datos.

En cambio, un codificador congelado conserva el conocimiento previo y reduce los
par\'ametros entrenables. Tambi\'en evita que un conjunto peque\~no degrade las
caracter\'isticas iniciales. Su principal limitaci\'on es que la pol\'itica solo puede
adaptar las capas posteriores a una representaci\'on fija.

Por \'ultimo, el ajuste fino permite adaptar los pesos preentrenados. Esta estrategia
ofrece mayor flexibilidad, aunque incrementa la memoria necesaria y puede producir
sobreajuste. La evidencia disponible no establece una regla universal: el resultado
cambia con la arquitectura, la tarea y el dominio de preentrenamiento
\cite{chi2025diffusion}, \cite{egbe2025dinov3dp}.

\section{S\'intesis y hueco de investigaci\'on}
\label{sec:hueco}

En conjunto, la literatura confirma que las pol\'iticas de difusi\'on representan
acciones multimodales y mantienen coherencia temporal. Los trabajos posteriores han
mejorado la geometr\'ia, la jerarqu\'ia y la velocidad de inferencia. Sin embargo, la
representaci\'on visual 2D suele permanecer fija o cambia junto con otros componentes.

De forma complementaria, los estudios sobre representaciones preentrenadas para
control descartan la existencia de un codificador universalmente superior. El
resultado depende del dominio, del objetivo de preentrenamiento y de la estrategia
de adaptaci\'on \cite{xiao2022mvp}, \cite{karamcheti2023voltron},
\cite{majumdar2023vc1}.

El art\'iculo original constituye una excepci\'on, puesto que compara ResNet-18,
ResNet-34 y CLIP ViT-B/16 bajo tres estrategias. No obstante, esa ablaci\'on se limita
a la tarea \textit{Square} \cite{chi2025diffusion}. En Push-T f\'isico solo se
contrastan el entrenamiento de extremo a extremo y dos codificadores congelados,
ImageNet y R3M.

Asimismo, DINOv3-DP ampl\'ia la evidencia a cuatro tareas e incluye Push-T. Ese
estudio no eval\'ua DINOv2 ni incorpora CLIP en la misma matriz comparativa
\cite{egbe2025dinov3dp}. Adem\'as, utiliza 100 \'epocas y una configuraci\'on de red
distinta de la implementaci\'on original empleada en este trabajo.

Los estudios generales de representaciones tampoco cierran este hueco, puesto que
emplean aprendizaje por refuerzo o clonaci\'on de comportamiento sobre bancos de
pruebas heterog\'eneos. Ninguno mantiene fija una \textit{Diffusion Policy} para
comparar estas alternativas bajo un protocolo com\'un de Push-T.

Por tanto, en el corpus revisado no se ha identificado una comparaci\'on conjunta de
cinco configuraciones concretas: ResNet-18 desde cero, ResNet-18 preentrenada en
ImageNet congelada y ajustada, DINOv2 congelado y CLIP congelado. Tampoco se informa
su relaci\'on entre rendimiento, par\'ametros entrenables, tiempo y memoria bajo un
protocolo \'unico de Push-T.

Este hueco justifica un estudio controlado que mantenga constante la pol\'itica
generativa, el conjunto de demostraciones y la evaluaci\'on. La variable experimental
queda as\'i restringida al bloque visual completo, entendido como la arquitectura del
codificador junto con sus pesos iniciales, su estrategia de adaptaci\'on y su cadena de
preprocesado, que en la pr\'actica var\'ian de forma conjunta. El \cref{cap:metodologia}
presenta el dise\~no empleado para realizar esa comparaci\'on y delimita hasta d\'onde
alcanza.
